\subsection{Nivel de abstracción 0: Diagrama de contexto}

  \paragraph{}En este diagrama se representará el funcionamiento de la
  aplicación de manera general, así como la interacción que mantiene el usuario
  con el sistema por medio de las entidades externas que generan flujos de
  entrada de información (Teclado, Ratón, Sistema y Unidad de almacenamiento) y
  los que reciben y procesan los flujos de salida de información del sistema
  (Pantalla, Impresora). A continuación se describirán las entidades externas
  que intervienen en el Diagrama de contexto.

  \begin{description}
   \item[Teclado] Permitirá al usuario introducir información en el sistema.

   \item[Ratón] Su función será la de permitir al usuario moverse por la
                aplicación.
   \item[Sistema] La información que proporcionará a la aplicación será la fecha
                  y hora del sistema en el que se encuentre instalada.
   \item[Unidad de almacenamiento] Representa toda unidad de almacenamiento
                                   externa que proporcionará archivos tales como
                                   imágenes, archivos pdf, sonidos, etc.
   \item[Servidor] Esta entidad interactúa con la aplicación proporcionando
                   flujos de información de entrada, cuando se le solicita
                   alguna página que tenga almacenada, y flujos de salida,
                   cuando la aplicación realice alguna operación que necesite
                   almacenar una página en el servidor.
   \item[Pantalla] Recibirá los flujos de información de salida de la aplicación
                   y será la encargada de mostrarlos al usuario.
   \item[Impresora] Recibirá información de salida de la aplicación y se
                    encargará de imprimirla en papel.
  \end{description}

  \paragraph{}La figura \ref{diagramaContexto} muestra el Diagrama de contexto.

        \begin{figure}[!ht]
            \begin{center}
            \includegraphics[]{08.Analisis_Funcional/8.2.DFDs/Niveles/Diagramas/diagrama_contexto.pdf}
            \caption{Diagrama de contexto.}
            \label{diagramaContexto}
            \end{center}
         \end{figure}